\begin{figure*}[htb]
\centering

% Styles
\tikzset{
    >=Stealth,
    font=\small,
    panelbox/.style={
        rounded corners,
        draw=black!60,
        very thick,
        inner sep=4mm,
        fill=white,
    },
    stage/.style={
        draw=black,
        thick,
        rounded corners,
        minimum width=3.2cm,
        minimum height=0.9cm,
        align=center,
        fill=black!3
    },
    mapbox/.style={
        draw=black,
        thick,
        rounded corners,
        minimum width=3.2cm,
        minimum height=0.9cm,
        align=center,
        fill=black!2
    },
}

\begin{tikzpicture}[
    node distance=12mm and 10mm
]

% --- Core nodes -------------------------------------------------------

% feature-space explanation a (left, top)
\node[stage, fill=black!2, minimum width=3.4cm] (anode) {
    concrete explanation\\[-1pt]
    $\displaystyle a \in \mathcal{A}$
};

% feature-space explanation a* (left, bottom; pushforward result)
\node[stage, fill=black!2, minimum width=3.4cm, below=5mm of anode] (astarnode) {
    tranferred explanation\\[-1pt]
    $\displaystyle a^{*} \in \mathcal{A}$
};

% pullback tau (middle, upper)
\node[mapbox, right=12mm of anode] (taunode) {
    pullback\\[-1pt]
    $\displaystyle \tau: \mathcal{A} \mapsto \mathcal{B}$
};

% pushforward psi (middle, lower)
\node[mapbox, right=12mm of astarnode] (psinode) {
    pushforward\\[-1pt]
    $\displaystyle \psi: \mathcal{B} \mapsto \mathcal{A}$
};

% problem-space explanation b* (right, top; pullback result)
\node[stage, fill=black!2, minimum width=3.4cm, right=12mm of taunode] (bstarnode) {
    transferred explanation\\[-1pt]
    $\displaystyle b^{*} \in \mathcal{B}$
};

% problem-space explanation b (right, bottom)
\node[stage, fill=black!2, minimum width=3.4cm, below=5mm of bstarnode] (bnode) {
    concrete explanation\\[-1pt]
    $\displaystyle b \in \mathcal{B}$
};

% Indicate closeness / correspondence within each space
\draw[<->, dashed, thick]
    (anode.south) -- node[right, font=\scriptsize] {$\approx$} (astarnode.north);

\draw[<->, dashed, thick]
    (bstarnode.south) -- node[right, font=\scriptsize] {$\approx$} (bnode.north);


% --- Arrows -----------------------------------------------------------

% Top row: a --tau--> b*
\draw[->, thick]
    (anode.east) -- (taunode.west);

\draw[->, thick]
    (taunode.east) -- (bstarnode.west);

% Bottom row: b --psi--> a*
\draw[->, thick]
    (bnode.west) -- (psinode.east);

\draw[->, thick]
    (psinode.west) -- (astarnode.east);

% Extra headroom at the top of each panel (no change to sides/bottom)
\coordinate (leftpad)  at ($(anode.north)+(0,7mm)$);
\coordinate (rightpad) at ($(bstarnode.north)+(0,7mm)$);

% --- Background region boxes (left/right only) -----------------------

\begin{scope}[on background layer]

\node[panelbox,
      fit=(anode) (astarnode) (leftpad),
      inner sep=2.5mm
] (leftpanel) {};

\node[panelbox,
      fit=(bstarnode) (bnode) (rightpad),
      inner sep=2.5mm
] (rightpanel) {};

\end{scope}

\node[font=\footnotesize,
      anchor=north,
      align=center]
    at ($(leftpanel.north)+(0,-1mm)$)
    {feature-space explanations\\$\mathcal{A} \subseteq \mathbb{R}^{|x|}$, e.g., $\mathcal{A} \subseteq \mathbb{R}^{2^{10}}$};

\node[font=\footnotesize,
      anchor=north,
      align=center]
    at ($(rightpanel.north)+(0,-1mm)$)
    {problem-space explanations\\$\mathcal{B} \subseteq \mathbb{R}^{|z|}$, e.g., $\mathcal{B} \subseteq \mathbb{R}^{2^{30}}$};

\end{tikzpicture}

\caption{Explanation Transformation. A feature-space explanation $a \in \mathcal{A}$ is pulled back via $\tau$ to yield a problem-space explanation $b^{*} \in \mathcal{B}$ close to $b$. Conversely, a problem-space explanation $b \in \mathcal{B}$ is pushed forward via $\psi$ to yield a feature-space explanation $a^{*} \in \mathcal{A}$ close to $a$.}
\label{fig:pullback-pushforward}
\end{figure*}
